\documentclass[10pt,letterpaper]{article}
\usepackage[spanish,es-noshorthands]{babel}
\usepackage[utf8]{inputenc}
\usepackage[T1]{fontenc}
\usepackage{lmodern}
\usepackage{amsmath,amssymb,booktabs,microtype}
\usepackage[margin=1.5cm,top=1.4cm,bottom=1.45cm]{geometry}
\usepackage{xcolor}
\usepackage{tikz,pgfplots}
\usepackage{multicol}
\usepackage{fancyhdr}
\usepackage{titlesec}
\usepackage{hyperref}
\pgfplotsset{compat=1.18}
\definecolor{ink}{HTML}{253342}
\definecolor{accent}{HTML}{8F465B}
\definecolor{muted}{HTML}{66727C}
\hypersetup{colorlinks=true,urlcolor=accent}
\titleformat{\section}{\normalfont\bfseries\color{accent}\normalsize}{\thesection.}{0.4em}{}
\titlespacing*{\section}{0pt}{5pt}{2pt}
\setlength{\parindent}{0pt}
\setlength{\parskip}{2pt}
\setlength{\columnsep}{0.7cm}
\pagestyle{fancy}
\fancyhf{}
\renewcommand{\headrulewidth}{0pt}
\fancyfoot[C]{\footnotesize\color{muted} Revista Internacional de Sentimientos Cuantificables \quad $\cdot$ \quad \thepage}
\begin{document}
\begin{center}
{\small\bfseries\color{accent} ARTÍCULO ORIGINAL \; | \; EDICIÓN ESPECIAL: 9 DE NOVIEMBRE}\par
\vspace{5pt}
{\LARGE\bfseries\color{ink} Sobre la imposibilidad estadística\\ de dejar de quererte}\par
\vspace{5pt}
{\large\itshape Un modelo de regresión lineal aplicado a 30 meses de amor por Cami}\par
\vspace{7pt}
{\normalsize Ing. Gómez, Esteban$^{1}$ \quad y \quad Camilla Echeverry$^{2}$}\par
{\small\color{muted} $^1$Departamento de Admiración Incondicional \quad $^2$Facultad de Estornudos Felinos}\par
{\footnotesize Recibido: hace dos años y medio \quad | \quad Aceptado: todos los días}\par
\end{center}
\vspace{3pt}
\noindent\colorbox{accent!7}{\begin{minipage}{\dimexpr\linewidth-2\fboxsep\relax}
\small\textbf{Resumen.} Este trabajo investiga la relación entre el tiempo compartido ($t$) y el cariño que siento por Cami ($A$). Se utiliza regresión lineal ordinaria por mínimos cuadrados sobre seis observaciones \emph{que me inventé}. El modelo presenta una pendiente positiva, aunque se reconoce una limitación fundamental: ningún método estadístico consigue medir con justicia lo tiernita que es mi mujer, ni lo orgulloso que estoy de verla cumplir sus sueños. \textbf{Palabras clave:} amor, mínimos cuadrados, futuro juntos, gatitos, cumpleaños.
\end{minipage}}
\setlength{\multicolsep}{6pt}
\small
\begin{multicols}{2}
\section{Introducción}
La hipótesis principal de este estudio es sencilla: desde que estamos juntos, mi vida se volvió más bonita. Después de dos años y medio, la evidencia cotidiana incluye abrazos, risas, apoyo mutuo y un fenómeno acústico especialmente adorable: los estornudos de Cami, prácticamente indistinguibles de los de un gatito.

A diferencia de otros estudios, la investigadora principal no sabe que está siendo objeto de esta publicación. Esto se considera necesario para preservar el factor sorpresa de su cumpleaños.

\section{Metodología}
Se define $t$ como los meses desde el inicio de nuestra relación y $A$ como un índice ficticio de cariño. Para obtener la recta de mejor ajuste se minimiza
\begin{equation}
\min_{\beta_0,\beta_1}\;\sum_{i=1}^{n}\bigl(A_i-\beta_0-\beta_1t_i\bigr)^2.
\end{equation}
De este modo, el estimador queda expresado como
\begin{equation}
\widehat A(t)=\widehat\beta_0+\widehat\beta_1t.
\end{equation}
Los datos de la Tabla~1 son \textbf{simbólicos}, creados para el chiste académico; no son mediciones reales de nuestros sentimientos.

\begin{center}
\small
\textbf{Tabla 1.} Observaciones ilustrativas\par\vspace{3pt}
\begin{tabular}{@{}rrrrrrr@{}}
\toprule
$t$ (meses)&0&6&12&18&24&30\\
\midrule
$A$ (índice)&2&3&4&6&7&9\\
\bottomrule
\end{tabular}
\end{center}

\section{Resultados}
La solución por mínimos cuadrados resulta en
\begin{equation}
\boxed{\widehat A(t)=1.667+0.2333t.}
\end{equation}
Así, a los 30 meses el modelo estima $\widehat A(30)\approx8.67$ unidades arbitrarias. La pendiente positiva respalda nuestra hipótesis ilustrativa: con el tiempo, crece lo que siento por ti.

\begin{center}
\begin{minipage}{\columnwidth}
\centering
\begin{tikzpicture}
\begin{axis}[
width=\columnwidth,height=3.4cm,
xlabel={Meses juntos},ylabel={Cariño (índice)},
xmin=-1,xmax=32,ymin=0,ymax=10,
xtick={0,6,12,18,24,30},ytick={0,2,4,6,8,10},
label style={font=\scriptsize},tick label style={font=\scriptsize},
grid=major,grid style={draw=gray!15},
axis line style={draw=muted},
]
\addplot[only marks,mark=*,mark size=1.65pt,color=accent] coordinates {(0,2)(6,3)(12,4)(18,6)(24,7)(30,9)};
\addplot[domain=0:30,samples=2,thick,color=ink] {1.6666667+0.2333333*x};
\end{axis}
\end{tikzpicture}
\par\footnotesize\color{muted} Figura 1. Datos ficticios y recta de ajuste.
\end{minipage}
\end{center}

\section{Discusión}
El modelo falla al intentar explicar tres variables decisivas: tu ternura, lo lindo que es verte sonreír y el orgullo inmenso que me dio verte graduarte como TCP. Sé que detrás de ese logro hubo dedicación, cansancio y muchísimo esfuerzo. Verte alcanzar esa meta me hizo admirarte todavía más.

También hay una variable que ninguna gráfica refleja: los dos nos la guerreamos trabajando para construir un futuro bien melo. A veces los días son difíciles, pero pensar que seguimos intentando, creciendo y soñando juntos me da fuerzas para continuar.

\section{Revisión por pares}
El comité felino emitió un estornudo y dio por terminada la sesión. El autor interpretó el dictamen como \emph{aceptado con abrazos} y se comprometió a entregar las correcciones en persona.

\section{Conclusión y dedicatoria}
Cami, feliz cumpleaños, mi amor. Entre tantas ecuaciones, hay algo que sí puedo afirmar sin margen de error: me encanta la personita que eres. Gracias por estos dos años y medio, por tu ternura, por tus estornuditos de gatito y por caminar conmigo hacia todo lo que soñamos.

Que este nuevo año te acerque a cada meta, a cada vuelo y a cada sueño que llevas en el corazón. Yo quiero seguir estando ahí, celebrándote en los días buenos y acompañándote en los difíciles.

\begin{equation*}
\boxed{\text{Mi lugar favorito}=\text{contigo}.}
\end{equation*}

\smallskip
\textbf{Trabajo futuro.} La próxima fase queda abierta: más cumpleaños, más vuelos que celebrar y más días normales de esos que se vuelven especiales porque los paso contigo.

\begin{flushright}
\textit{Te amo muchísimo.}\\
\textbf{Tu novio $\heartsuit$}
\end{flushright}
\end{multicols}
\end{document}
